\documentclass[11pt]{scrartcl}
\usepackage[utf8]{inputenc}
\usepackage{evank}


%% insert title here
\newcommand{\titlee}{Modern Physics [Solutions]}
\parindent 0 pt

\newcommand{\officialsol}[1]{See the official solutions \href{#1}{here}.}

\title{\titlee}
\author{Cambridge Physics Academy}
\date{2023-2024}


\begin{document}

\thispagestyle{empty}
\begin{center}
    
    
    \Huge \textbf{\textsf{\titlee}}

    \large \theauthor
\end{center}

\setcounter{section}{-1}
\section{Readings}
\begin{itemize}
    \item Ch 45-50 of HRK
\end{itemize}
Feel free to do problems from the readings as extra practice.


\section{Lecture review}

Quantum mechanics is the extension of wave-particle duality to all matter! Everything becomes probabilistic. Some fundamental ideas are:
\begin{itemize}
    \item We have probabilities, not determined outcomes
    \item We often find quantized phenomena: energy levels, flux, etc
    \item Only operates at the small scale. More rigor on this in just a second.
\end{itemize}

\begin{defn}[de-Broglie Relations]
    The de-Broglie wavelength gives a length scale for when quantum mechanical effects matter. It is defined as
    $$\lambda = \frac hp \iff p = \hbar k.$$
    There is also a de-Broglie frequency,
    $$f = \frac Eh \iff E = \hbar \omega.$$
\end{defn}


\begin{theorem}[Heisenberg Uncertainty Principle]
    There are several pairs of quantities which have an uncertainty principle, but the two most famous are $x \leftrightarrow p$ and $E \leftrightarrow t$. Mathematically, it states
    $$\Delta x \Delta p \ge \frac{\hbar} 2, \Delta E \Delta t \ge \frac \hbar 2.$$
    In most problems, the numerical factors will not matter.
\end{theorem}

\begin{idea}[WKB Approximatino]
    In the regime where the debroglie wavelength is much smaller than the length scale of the problem, we can treat particles semi-classically! Using this we have that for a quantum ``standing wave'',
    $$\oint k dx = 2\pi n \implies \oint \frac{p}{\hbar} dx = 2\pi n.$$
    This reduces to
    $$\oint p dx = nh.$$
\end{idea}

Most other ``modern physics'' olympiad questions are just questions that apply a lot of the old concepts that we’ve learned so far. The most popular topics to use are Relativistic Dynamics, Electrodynamics, and Gravity.

\section{Problem Set}

\subsection{Exercises}

\begin{prob}
    Use the WKB approximation to find the quantization condition for a particle in a box.
\end{prob}
\begin{sol}
    Since for a particle in a box there is no potential within the box, we have $p = \sqrt{2mE}$. Thus, 
    $$\oint p\, dx = \boxed{2L\sqrt{2mE} = nh}.$$
\end{sol}

\begin{prob}
    An atom of mass $m$ with velocity $v$ travels toward a photon of frequency $\omega$. The atom absorbs are re-emits the photon in the opposite direction. What is the new velocity of the mass?
\end{prob}
\begin{sol}
    In the frame of the atom, the photon is blue shifted. Then, as it emits it, the photon gets further blue shifted. So, the final frequency of the photon is
    $$\omega' = \frac{1 + v/c}{1-v/c}\omega.$$
    Thus, by conservation of momentum,
    $$mv - \frac{\hbar \omega}{c} = mv' + \frac{\hbar \omega'}{c}.$$
    So, solving for $v'$, we get
    $$v' = v - \frac{\hbar\omega}{c}\frac{2}{1-v/c}.$$
\end{sol}

\begin{prob}
    \href{https://aapt.org/physicsteam/2015/upload/E3-2-5.pdf#page=3}{USAPhO 2015/A1}.
\end{prob}
\begin{sol}
    \officialsol{https://aapt.org/physicsteam/upload/USAPhO-2015-Solutions.pdf}
\end{sol}

\begin{prob}
    \href{https://aapt.org/physicsteam/2019/upload/USAPhO-2018.pdf#page=8}{USAPhO 2018/B2}
\end{prob}
\begin{sol}
    \officialsol{https://aapt.org/physicsteam/upload/USAPhO-2018-Solutions.pdf}
\end{sol}

The following problem is out of scope because it uses linear algebra, but if you do know linear algebra, it'll be useful for getting practice with a problem on an unfamiliar topic -- sort of the theme of all the ``modern physics'' olympiad problems.

\begin{prob}
    Electrons are either spin up or spin down in the $z$-direction, but before they are measured, they can be in a superposition of the two states. Because those two states are opposites, we can actually think of them as an orthogonal basis of a 2 dimensional vector space. Thus the quantum state representing the spin of any electron can be written as
    $$\ket{\psi} = c_0 \binom 10 + c_1 \binom 01,$$
    where $c_0,c_1 \in \mathbb C$, and $|c_0|^2 + |c_1|^2 = 1$. The fact that they sum to $1$ suggest that $|c_0|^2$ is the probability that it is $\langle 1,0\rangle$ (spin up) and $|c_1|^2$ is the probability that it is $\langle 0,1\rangle$ (spin down). But there are more bases than spin up and down in the $z$-direction. We can also have spin up and down in the $x$ direction:
    $$\ket{x+} = \frac{1}{\sqrt{2}}\binom 11, \ket{x-} = \frac{1}{\sqrt{2}} \binom{1}{-1}.$$
    \begin{anum}
        \item Suppose you prepare a quantum state $\ket \psi = \ket{x+}$. You then measure its spin in the $z$-direction. What is the probability that it is spin up?
        \item Now suppose you prepare a quantum state $\ket \psi = \langle 1,0\rangle$. You then measure it in the $x$-direction. What is the probability that it is spin up in the $x$ direction?
        \item You then put your spin up electron under some sort of magnetic field such that it begins evolving through time. This magnetic field is such that when an electron is in the state $\ket{x+}$, it has energy $+g$ and when it is in the state $\ket{x-}$, it has energy $-g$. What is the average energy of a $\ket{z+}$ electron?
        \item After solving the schrodinger equation, you find that as a function of time,
        $$\ket{\psi(t)} = \begin{pmatrix}\cos \left(\frac{gt}{2\hbar}\right) \\ \sin \left(\frac{gt}{2\hbar}\right)\end{pmatrix}.$$
        Show that this is consistent with Heisenberg's uncertainty principle by estimating $\Delta E$ and $\Delta t$.
    \end{anum}
\end{prob}
\begin{alphasol}
    \item We can write $\ket{x+}$ as a linear combination of $z$ states,
    $$\ket{x+} = \frac{1}{\sqrt 2} \ket{z+} + \frac{1}{\sqrt 2}\ket{z-},$$
    so the probability that it is spin up is $\boxed{1/2}$ because it is an equal combination of the two.
    \item We do the same strategy as in the previous part,
    $$\langle 1,0\rangle = \frac{1}{\sqrt 2}\ket{x+} + \frac{1}{\sqrt 2}\ket{x-}.$$
    \item As discussed in part (a) the $\ket{z+}$ state is an equal superposition of $\ket{x+}$ and $\ket{x-}$, so
    $$E_{\text{avg}} = \frac 12 (+g) + \frac 12 (-g) = 0.$$
    \item The time for the state to return to the same state is  period $T = \frac{2\pi}{g/2\hbar} \sim \frac{\hbar}{g}$. Now, the energy goes from $+g$ to $-g$, so $\Delta E \sim g$, so $$\Delta E \Delta t \sim g \cdot \frac{\hbar}{g} \sim \hbar,$$
    consistent with Heisenberg's uncertainty principle.
\end{alphasol}


The below two problems are both from the IPhO, but they should both be pretty approachable.

\begin{prob}
    \href{https://ipho.olimpicos.net/pdf/IPhO_2006_Q1.pdf}{IPhO 2006/T1}. 
\end{prob}
\begin{sol}
    \officialsol{https://ipho.olimpicos.net/pdf/IPhO_2006_S1.pdf}
\end{sol}

\begin{prob}
    \href{https://ipho.olimpicos.net/pdf/IPhO_2009_Q2.pdf}{IPhO 2009/T2}
\end{prob}
\begin{sol}
    \officialsol{https://ipho.olimpicos.net/pdf/IPhO_2009_S2.pdf}
\end{sol}

\begin{prob}
    \href{https://olympiads.hbcse.tifr.res.in/wp-content/uploads/2020/02/INPHO2020-Questions-en.pdf}{INPhO 2020/3}. This question looks simple, but be careful.
\end{prob}
\begin{sol}
    \officialsol{https://olympiads.hbcse.tifr.res.in/wp-content/uploads/2020/03/INPHO2020-Solutions-20200220.pdf}
\end{sol}


\subsection{Challenge Problems}

\begin{prob}
    \href{https://aapt.org/Common2022/upload/2021-USAPhO_plus.pdf#page=5}{USAPhO+ 2021/T2}
\end{prob}
\begin{sol}
    \officialsol{https://aapt.org/Common2022/upload/2021-USAPhO_plus_Solutions_v2.pdf}
\end{sol}

\begin{prob}
    \href{https://aapt.org/Common2022/upload/2021-USAPhO_plus.pdf#page=8}{USAPhO+ 2021/T3}
\end{prob}
\begin{sol}
    \officialsol{https://aapt.org/Common2022/upload/2021-USAPhO_plus_Solutions_v2.pdf}
\end{sol}

\begin{prob}
    \href{https://apho.olimpicos.net/pdf/APhO_2016_Q3.pdf}{APhO 2016/T3}.
\end{prob}
\begin{sol}
    \officialsol{https://apho.olimpicos.net/pdf/APhO_2016_S3.pdf}
\end{sol}

\begin{prob}
    \href{https://aapt.org/physicsteam/upload/2022_Theory_TST.pdf#page=5}{US Physics Team 2022 TST Q3}, a problem on gravitational waves.
\end{prob}
\begin{sol}
    \officialsol{https://aapt.org/physicsteam/upload/2022_Theory_TST_Solutions.pdf}
\end{sol}

\subsection{USAPhO Practice}

\begin{prob}[USAPhO 2022/B3]
    The kinetic energy $E$, momentum $p$, and velocity $v$ of a particle moving in one dimension satisfy
    $$F = \frac{dp}{dt} = \frac{dE}{dx}, v = \frac{dE}{dp}$$
    where $F$ is the external force. For a free particle, the momentum and energy are related by $E = p^2/2m$. However, when an electron moves inside a metal, its interactions with the crystal
lattice of positively charged ions lead to a different relationship between momentum and energy. All
of the above identities still apply, but now suppose that
$$E(p) = V(1-\cos (pb))$$
where $V$ and $b$ are constants that depend on the metal. This result is inherently quantum mechanical
in origin, and as we will see, it leads to some rather strange behavior.
\begin{anum}
    \item First, we investigate the motion of the electron in general.
    \begin{enumerate}[label=\roman*.]
        \item Find the velocity as a function of $p$.
        \item The effective mass $m_*$ of the electron is defined so that it satisfies $F = m_* a$. Find the effective
mass as a function of $p$.
    \end{enumerate}
    \item Now suppose a metal rod of infinite length, aligned with the $x$-axis, contains conducting electrons
of charge $-e$, initially with zero momentum. At time $t = 0$, an electric field $\mathbf E = E_0 \hat{\mathbf x}$ is turned
on, and is experienced by every electron in the rod. Ignore the interactions of the electrons with
each other.
\begin{enumerate}[label=\roman*.]
    \item For an electron that starts at $x = 0$ at time $t = 0$, find its position $x(t)$.
    \item If the number of conducting electrons per unit volume is $n$, and the cross-sectional area of the
rod is $A$, find the average current in the rod over a long time.
    \item Now suppose that every time $\tau$, each electron suffers a collision with the crystal lattice, causing
its momentum to reset to zero. In the limit of frequent collisions, $eE_0 b \tau \ll 1$, find the average
current in the rod over a long time.
\item If $\tau$ can be freely adjusted, estimate the maximum possible average current in the rod, up to
a dimensionless constant.
\end{enumerate}
\end{anum}
\end{prob}
\begin{sol}
    \officialsol{https://aapt.org/physicsteam/upload/2022-USAPhO-solutions.pdf}
\end{sol}

\begin{prob}[USAPhO 2011/A4]
    In this problem we consider a simplified model of the electromagnetic radiation inside a cubical
box of side length $L$. In this model, the electric field has spatial dependence
$$E(x, y, z) = E_0 \sin(k_xx) \sin(k_yy) \sin(k_zz).$$
where one corner of the box lies at the origin and the box is aligned with the $x$, $y$, and $z$ axes. Let $h$ be Planck’s constant, $k_B$ be Boltzmann’s constant, and $c$ be the speed of light.
\begin{anum}
    \item The electric field must be zero everywhere at the sides of the box. What condition does this
impose on $k_x$, $k_y$, and $k_z$? (Assume that any of these may be negative, and include cases
where one or more of the $k_i$
is zero, even though this causes $E$ to be zero.)
    \item In the model, each permitted value of the triple $(k_x, k_y, k_z)$ corresponds to a quantum state.
These states can be visualized in a \textit{state space}, which is a notional three-dimensional space
with axes corresponding to $k_x$, $k_y$, and $k_z$. How many states occupy a volume $s$ of state space,
if $s$ is large enough that the discreteness of the states can be ignored?
\item Each quantum state, in turn, may be occupied by photons with frequency $\omega = \frac f{2\pi} = c|\mathbf k |$, where
$$|\mathbf k| = \sqrt{k_x^2 + k_y^2+k_z^2}.$$
In the model, if the temperature inside the box is $T$, no photon may have energy greater than $k_BT$. What is the shape of the region in state space corresponding to occupied states?
\item As a final approximation, assume that each occupied state contains exactly one photon. What
is the total energy of the photons in the box, in terms of $h, k_B, c, T$, and the volume of the box
$V$ ? Again, assume that the temperature is high enough that there are a very large number of
occupied states. (Hint: divide state space into thin regions corresponding to photons of the
same energy.)
\end{anum}
\end{prob}
\begin{sol}
    \officialsol{https://aapt.org/Common/upload/USAPhO-2011-Solutions_rev-4012021.pdf}
\end{sol}


\end{document}